\section{结论}\label{sec:conc}

本文提出了一种用于实现递推条件概率滤波的新算法，即条件自举粒子滤波器。该方法以一组随机样本直接表示所需的后验状态分布，从而彻底摆脱了解析可处理性假设对系统与观测模型形式的束缚。

然而，在诸多具有实际意义的工程问题中，样本集内互异状态值的数量可能在若干步迭代后迅速塌缩。为缓解这一样本贫化难题，本文提出了两项带有一定启发式特征的补救方案——样本抖动加糙过程与单步超前先验筛选方案，并在具有四维状态空间的非线性纯方位目标追踪仿真中予以验证。仿真结果表明，即便在样本规模相当适中的条件下，条件自举滤波器的估计精度与不确定度校准能力仍远优于标准一阶线性化正态滤波器。

后续仍需开展进一步研究，以建立具有严格理论基础的实用机制来防止样本多样性的退化塌缩；同时，制定关于为获得满意自举滤波性能所需样本规模的定量选取指南也将极具实用价值。此外，还应开展更为系统的随机抽样仿真（原称 Monte Carlo simulation）评估，以便在各类重要非线性状态估计问题中定量刻画滤波性能。

\section{参考文献}\label{sec:refs}

\bibliography{html/src/refs}

\htdfntext{1}{本文原刊发表于 \emph{IEE Proceedings-F (Radar and Signal Processing)}，第 140 卷第 2 期，第 107–113 页，1993 年 4 月。论文编号 Paper 9241F (E5)，首次收稿日期为 1992 年 4 月 27 日，修订稿收稿日期为 1992 年 10 月 8 日。}

\htdfntext{2}{N. J. Gordon 与 D. J. Salmond 撰写本文时任职于英国 Farnborough 的 Defence Research Agency；A. F. M. Smith 任职于英国伦敦 Imperial College 统计系。本研究工作得到 DRA Farnborough 资助支持。}

% <!--FOOTNOTES-->
